\kchapter{Volumes and Persistent Storage}{pvc}

\begin{objectives}
\begin{itemize}
  \item why container filesystems are temporary;
  \item ephemeral volumes: \kc{emptyDir}, \kc{configMap}, \kc{secret}, \kc{projected};
  \item PersistentVolumes, PersistentVolumeClaims, and StorageClasses;
  \item access modes, reclaim policies, volume expansion, and snapshots;
  \item the CSI driver model.
\end{itemize}
\end{objectives}

\section{The problem}

A container's filesystem lives and dies with the container. Restart the
container and every file it wrote is gone. Databases, uploads, and caches need
storage with a different lifetime. Kubernetes offers \kterm{volumes}: directories
mounted into containers whose lifetime is managed separately.

\section{Ephemeral volumes}

These live as long as the \emph{Pod} (surviving container restarts, but not Pod deletion).

\begin{center}
\ktablestart
\begin{tabularx}{\linewidth}{P{3.0cm} L}
\kheader \kth{Volume} & \kth{Use} \\
\kc{emptyDir} & Scratch space shared between containers in a Pod. Set
                \kc{medium: Memory} for a RAM disk. \\
\kc{configMap}, \kc{secret} & Expose config or credentials as files. \\
\kc{projected} & Combine several sources (configMap, secret, downwardAPI,
                 serviceAccountToken) into one directory. \\
\kc{downwardAPI} & Expose Pod labels, annotations, and resources as files. \\
Generic ephemeral & A per-Pod PVC that is deleted with the Pod. \\
\end{tabularx}
\end{center}

\begin{warningbox}
\kc{hostPath} mounts a directory from the node. It ties the Pod to one node and
can give a container access to the host --- a serious security risk. Avoid it
except for system agents that truly need it.
\end{warningbox}

\section{Persistent storage: three objects}

\begin{center}
\resizebox{\linewidth}{!}{%
\begin{tikzpicture}[
  o/.style={draw=k8sblue, fill=white, line width=1pt, rounded corners=5pt,
            minimum width=3.2cm, minimum height=1.9cm, align=center, font=\small\bfseries,
            text=k8snavy},
  >={Stealth[length=2.6mm]}]
  \node[o] (pod) at (0,0) {\kimg[0.8cm]{pod}\\Pod\\[-2pt]\scriptsize\mdseries mounts the claim};
  \node[o] (pvc) at (4.4,0) {\kimg[0.8cm]{pvc}\\PersistentVolume-\\Claim\\[-2pt]\scriptsize\mdseries ``I need 10Gi RWO''};
  \node[o] (pv)  at (8.8,0) {\kimg[0.8cm]{pv}\\PersistentVolume\\[-2pt]\scriptsize\mdseries a real 10Gi disk};
  \node[o] (sc)  at (8.8,-3.2) {\kimg[0.8cm]{sc}\\StorageClass\\[-2pt]\scriptsize\mdseries ``fast-ssd'' recipe};
  \node[o, draw=k8sgray] (disk) at (13.2,0) {\Huge\faHdd\\[2pt]Cloud disk / NFS\\[-2pt]\scriptsize\mdseries EBS, PD, Azure Disk\ldots};
  \draw[->, thick, k8sblue] (pod) -- node[above, font=\scriptsize]{uses} (pvc);
  \draw[<->, thick, k8sblue] (pvc) -- node[above, font=\scriptsize]{bound} (pv);
  \draw[->, thick, k8sorange, dashed] (sc) -- node[right, font=\scriptsize, align=left]{dynamically\\provisions} (pv);
  \draw[->, thick, k8sgray] (pv) -- node[above, font=\scriptsize]{represents} (disk);
  \draw[->, thick, k8sorange, dashed] (pvc.south) |- node[pos=0.75, below, font=\scriptsize]{storageClassName} (sc.west);
\end{tikzpicture}}
\end{center}

\begin{itemize}
  \item A \kicon{pv}~\kterm{PersistentVolume} (PV) is a piece of storage in the
        cluster --- cluster-scoped, like a node is a piece of compute.
  \item A \kicon{pvc}~\kterm{PersistentVolumeClaim} (PVC) is a namespaced
        \emph{request} for storage: size, access mode, class.
  \item A \kicon{sc}~\kterm{StorageClass} describes a \emph{kind} of storage and
        which driver creates it, enabling \kterm{dynamic provisioning}: when a
        PVC appears, a matching PV (and real disk) is created automatically.
\end{itemize}

\begin{analogy}
A StorageClass is a menu item (``large SSD''), a PVC is your order, and the PV
is the plate that arrives. You never walk into the kitchen yourself.
\end{analogy}

\begin{yamlcode}{storageclass.yaml}
apiVersion: storage.k8s.io/v1
kind: StorageClass
metadata:
  name: fast-ssd
  annotations:
    storageclass.kubernetes.io/is-default-class: "false"
provisioner: ebs.csi.aws.com        # the CSI driver (cloud specific)
parameters:
  type: gp3
reclaimPolicy: Delete               # or Retain
allowVolumeExpansion: true
volumeBindingMode: WaitForFirstConsumer   # create the disk in the zone where the Pod lands
\end{yamlcode}

\begin{yamlcode}{pvc-and-pod.yaml}
apiVersion: v1
kind: PersistentVolumeClaim
metadata:
  name: uploads
spec:
  accessModes: ["ReadWriteOnce"]
  storageClassName: fast-ssd        # omit to use the cluster default class
  resources:
    requests:
      storage: 10Gi
---
apiVersion: v1
kind: Pod
metadata:
  name: uploader
spec:
  containers:
    - name: app
      image: nginx:1.27
      volumeMounts:
        - name: data
          mountPath: /usr/share/nginx/html
  volumes:
    - name: data
      persistentVolumeClaim:
        claimName: uploads
\end{yamlcode}

\begin{bashcode}{terminal}
kubectl get storageclass           # kind ships "standard" (local-path) as default
kubectl apply -f pvc-and-pod.yaml
kubectl get pvc,pv
kubectl describe pvc uploads       # Pending? read the events
\end{bashcode}

\section{Access modes}

\begin{center}
\ktablestart
\begin{tabularx}{\linewidth}{P{3.6cm} P{1.6cm} L}
\kheader \kth{Mode} & \kth{Short} & \kth{Meaning} \\
\kc{ReadWriteOnce} & RWO & Mounted read-write by one \emph{node}. Block disks (EBS, PD). Most common. \\
\kc{ReadWriteOncePod} & RWOP & Read-write by exactly one \emph{Pod}. Strongest guarantee. \\
\kc{ReadOnlyMany} & ROX & Read-only by many nodes. \\
\kc{ReadWriteMany} & RWX & Read-write by many nodes. Needs a shared filesystem (NFS, EFS, CephFS, Azure Files). \\
\end{tabularx}
\end{center}

\section{Lifecycle and reclaim policy}

When a PVC is deleted, the bound PV's \kc{reclaimPolicy} decides the fate of the data:

\begin{itemize}
  \item \kc{Delete} --- the PV and underlying disk are deleted (default for dynamic provisioning);
  \item \kc{Retain} --- the PV and disk are kept for manual recovery. Use it for anything precious.
\end{itemize}

\begin{warningbox}
With \kc{reclaimPolicy: Delete}, \kc{kubectl delete pvc} or deleting a
namespace \textbf{destroys the data}. For production databases use a
\kc{Retain} class and take backups or snapshots.
\end{warningbox}

\section{Expansion and snapshots}

If the StorageClass has \kc{allowVolumeExpansion: true}, grow a volume by editing
the PVC's \kc{spec.resources.requests.storage}. Shrinking is not supported.

With a CSI driver that supports it, take point-in-time snapshots:

\begin{yamlcode}{snapshot.yaml}
apiVersion: snapshot.storage.k8s.io/v1
kind: VolumeSnapshot
metadata:
  name: uploads-2026-10-10
spec:
  volumeSnapshotClassName: csi-snapclass
  source:
    persistentVolumeClaimName: uploads
---
# Restore: a new PVC whose dataSource is the snapshot
apiVersion: v1
kind: PersistentVolumeClaim
metadata:
  name: uploads-restored
spec:
  accessModes: ["ReadWriteOnce"]
  resources:
    requests:
      storage: 10Gi
  dataSource:
    name: uploads-2026-10-10
    kind: VolumeSnapshot
    apiGroup: snapshot.storage.k8s.io
\end{yamlcode}

\section{CSI: the storage plugin model}

The \kterm{Container Storage Interface} lets storage vendors ship drivers
independently of Kubernetes releases. Every modern backend --- AWS EBS/EFS, GCP
PD/Filestore, Azure Disk/Files, Ceph (Rook), Longhorn, Portworx, NetApp --- has a
CSI driver. You install the driver, create a StorageClass that names it, and
PVCs do the rest.

\begin{recap}
\begin{itemize}
  \item Container filesystems are temporary; volumes give data a separate lifetime.
  \item PVC = request, PV = actual storage, StorageClass = recipe for dynamic provisioning.
  \item Choose access modes and reclaim policies deliberately; \kc{Delete} really deletes.
  \item Use snapshots and backups for anything you cannot lose.
\end{itemize}
\end{recap}
